\section{Feasibility of a Learned Transition Solver}

\subsection{Problem Statement}

The current project goal is not to replace a full high-fidelity hydrodynamic model.
Instead, the learned model is expected to serve as a short-horizon state transition solver
inside a controller, a lightweight simulator, or a reinforcement-learning environment.

Let the latent engineering state be
\[
\mathbf{x}_k =
\begin{bmatrix}
\mathbf{a}_k \\
\boldsymbol{\omega}_k \\
\mathbf{v}_k
\end{bmatrix}
\in \mathbb{R}^{9},
\]
where $\mathbf{a}_k$ is body-frame linear acceleration,
$\boldsymbol{\omega}_k$ is body-frame angular rate,
and $\mathbf{v}_k$ is body-frame velocity.

The control-related input is
\[
\mathbf{u}_k \in \mathbb{R}^{8},
\]
and the auxiliary context is
\[
\mathbf{c}_k =
\begin{bmatrix}
\phi_k & \theta_k & \sin\psi_k & \cos\psi_k & \mathbf{p}_k
\end{bmatrix}^{\top},
\]
where $\phi_k,\theta_k,\psi_k$ are attitude terms and $\mathbf{p}_k$ denotes motor power observations.

The desired one-step transition operator is
\[
\mathbf{x}_{k+1} = f_{\vartheta}(\mathbf{x}_k, \mathbf{u}_k, \mathbf{c}_k),
\]
or equivalently in increment form
\[
\Delta \mathbf{x}_k = g_{\vartheta}(\mathbf{x}_k, \mathbf{u}_k, \mathbf{c}_k),
\qquad
\mathbf{x}_{k+1} = \mathbf{x}_k + \Delta \mathbf{x}_k.
\]

\subsection{Why Sensor Preprocessing is Necessary}

Raw asynchronous sensor streams do not directly satisfy the assumptions required by a transition model:

\begin{enumerate}
\item The observations arrive at different rates.
\item DVL velocity is sparse and intermittent.
\item IMU signals are dense but noisy and bias-sensitive.
\item Power observations are low-rate auxiliary measurements rather than primary state labels.
\end{enumerate}

If these signals are fed directly into a neural predictor, the network must simultaneously solve:

\begin{enumerate}
\item time alignment,
\item denoising and bias suppression,
\item state estimation,
\item and true controlled dynamics learning.
\end{enumerate}

This overloads the model and usually leads to unstable rollout.

\subsection{Causal Filtering Makes the Learning Problem Better Posed}

Suppose a causal estimator produces a proxy state
\[
\hat{\mathbf{x}}_k = \mathcal{F}( \mathcal{Y}_{0:k}, \mathbf{u}_{0:k}),
\]
where $\mathcal{Y}_{0:k}$ contains all sensor observations available up to time $k$.

The key requirement is causality:
\[
\hat{\mathbf{x}}_k = \mathcal{F}(\mathcal{Y}_{0:k}, \mathbf{u}_{0:k}),
\]
that is, the estimate at time $k$ must not depend on future observations $\mathcal{Y}_{k+1:T}$.
This requirement is essential because the downstream transition model is intended for online rollout.

Under this setup, the learning target becomes
\[
\hat{\mathbf{x}}_{k+1} \approx f_{\vartheta}(\hat{\mathbf{x}}_k, \mathbf{u}_k, \mathbf{c}_k),
\]
which is a standard output-error system identification problem with filtered state proxies.

This is mathematically feasible as long as the filtered proxy states satisfy two properties:

\begin{enumerate}
\item They are sufficiently low-noise to reduce label variance.
\item They preserve the dominant short-horizon transition structure relevant to control.
\end{enumerate}

The existing KF / ESKF preprocessing chain is therefore not a cosmetic improvement.
It reduces the estimation burden seen by the neural transition model.

\subsection{Why Proxy-State Training Is More Representative}

For a transition solver, the decisive property is not whether the training labels preserve every sensor-side detail,
but whether they preserve the controlled short-horizon evolution relevant to recursive prediction.

The proxy-state formulation is more representative than raw-sensor fitting because:
\begin{enumerate}
\item it places both input and output on the same aligned time base;
\item it expresses the learned quantity in a body-frame state space compatible with recursive rollout;
\item it separates ``state construction'' from ``transition learning'';
\item it allows reliability and freshness information to enter the model explicitly through context variables.
\end{enumerate}

Hence, the learned predictor is no longer asked to imitate a heterogeneous observation stream.
It is asked to approximate a state transition law on a causally estimated state manifold.

\subsection{Role of DVL, IMU, and Power}

The sensors have different mathematical roles:

\begin{itemize}
\item IMU provides dense propagation information.
\item DVL provides sparse but strong velocity correction.
\item Power provides auxiliary evidence about actuator loading and effective thrust.
\end{itemize}

Hence, the engineering state should not be equated with any single sensor.
Instead, it should be regarded as a causally estimated proxy state.

This also explains why power is better treated as context or auxiliary observation rather than as a direct state target:
power helps constrain the map from command to effective actuation, but it is not itself the rigid-body motion state.

\subsection{Why Shared Scaling is a Hard Requirement}

Assume the rollout is trained in normalized space:
\[
\tilde{\mathbf{x}} = \frac{\mathbf{x} - \boldsymbol{\mu}}{\boldsymbol{\sigma}}.
\]

If the initial state $\mathbf{x}_k$ is normalized by one scaler
and the target $\mathbf{x}_{k+1}$ is normalized by another scaler,
then the recurrence
\[
\tilde{\mathbf{x}}_{k+1} = \tilde{\mathbf{x}}_k + \Delta \tilde{\mathbf{x}}_k
\]
no longer corresponds to a consistent physical-space update.

Therefore, for every shared state component used both in the history input and in the prediction target,
the same normalization statistics must be used.

This is not a minor implementation detail.
It is a necessary condition for the increment-based rollout semantics to remain valid.

\subsection{Phase 1 Feasibility Claim}

After enforcing:

\begin{enumerate}
\item causal filtered state proxies, and
\item shared normalization semantics for overlapping state dimensions,
\end{enumerate}

the project has a defensible foundation for learning a short-horizon transition solver.

What Phase 1 proves is:

\begin{itemize}
\item the training target can be interpreted as a causal state-proxy transition problem;
\item the rollout objective is numerically coherent;
\item the model can be upgraded toward a one-step solver without rebuilding the whole data chain.
\end{itemize}

What Phase 1 does \emph{not} prove is:

\begin{itemize}
\item global closed-loop stability,
\item optimal controller compatibility,
\item or equivalence to a high-fidelity physical simulator.
\end{itemize}

\subsection{Next Mathematical Step}

The next step is to shift from block prediction toward explicit one-step transition learning:
\[
\hat{\mathbf{x}}_{k+1} = f_{\vartheta}(\hat{\mathbf{x}}_k, \mathbf{u}_k, \mathbf{c}_k),
\]
with short recursive unrolling during training.

The corresponding loss should combine:
\[
\mathcal{L}
=
\lambda_1 \mathcal{L}_{\text{step}}
+
\lambda_2 \mathcal{L}_{\text{rollout}}
+
\lambda_3 \mathcal{L}_{\text{obs}}
+
\lambda_4 \mathcal{L}_{\text{stability}},
\]
where:
\begin{itemize}
\item $\mathcal{L}_{\text{step}}$ enforces one-step transition accuracy,
\item $\mathcal{L}_{\text{rollout}}$ enforces short recursive consistency,
\item $\mathcal{L}_{\text{obs}}$ enforces agreement with sensor-space observations,
\item $\mathcal{L}_{\text{stability}}$ penalizes divergence, non-finite behavior, or excessive tail growth.
\end{itemize}

This is the mathematically natural bridge from offline fitted sequence prediction
to a transition model usable by control and reinforcement learning.
